\section{PAPER 028: Studio v2 — A Six-Panel Live Dashboard for Autonomous Knowledge Graph Operations}

\textbf{Author}: Bryan Alexander Buchorn  
\textbf{Affiliation}: Independent Researcher, Las Vegas, NV, USA  
\textbf{Status}: v2.52.0 (Phase 172 (STRB) COMPLETE)
\textbf{Date}: May 2, 2026

\textbf{CEREBRUM Phases 75 + 78}

---

\subsection{Abstract}

We present \textbf{Studio v2}, an extension of CEREBRUM's Glass-Box Reasoning Studio (Phase 20) with six live monitoring panels designed for autonomous discovery operations. Studio v2 introduces an attachment API (\texttt{attach\_research\_agent}, \texttt{attach\_modulator}, \texttt{attach\_loop}, \texttt{attach\_provenance\_ledger}) that wires optional engines to the dashboard without hard depen\-dencies. The six panels cover: (1) AutoApprover audit log, (2) Contra\-dictionResolver revision queue, (3) DiscoveryCalibrator community heatmap, (4) ChemicalModulator blood panel, (5) Autonomous Loop cycle history, and (6) Provenance Panel (Phase 78) — a three-part graph provenance view showing batch bar chart and cycle timeline. All panels degrade gracefully when the corre\-sponding engine is not attached, returning placeholder output rather than raising exceptions. This design allows progressive adoption: operators attach only the engines relevant to their deployment.

---

\subsection{1. Motivation: Observ\-ability for Autonomous Systems}

The original Reasoning Studio [Buchorn, 2026] was designed for human-in-the-loop query forensics: a user submits a query, inspects the reasoning path, and adjusts parameters. Studio v2 addresses a different use case: \textbf{long-running autonomous discovery operations} where an operator needs continuous visibility into the health of the full discovery pipeline without requiring per-query interaction.

Key obser\-vability requi\-rements:
\begin{itemize}
\item \textbf{Decision trans\-parency}: What is AutoApprover deciding, and why?
\item \textbf{Quality monitoring}: Is the circuit breaker at risk of tripping?
\item \textbf{Coverage monitoring}: Which communities are being over- / under-explored?
\item \textbf{Physio\-logical state}: Is ChemicalModulator operating near homeostatic baseline?
\item \textbf{Audit trail}: What was mater\-ialized, when, and has anything been rolled back?

\end{itemize}
---

\subsection{2. Attachment API}

\begin{lstlisting}
studio = StudioEngine(graph)

studio.attach_research_agent(agent)
studio.attach_modulator(modulator)
studio.attach_loop(loop)
studio.attach_provenance_ledger(ledger)
\end{lstlisting}

Each attachment is independent. A \texttt{StudioEngine} with no attachments still renders the original Phase 20 reasoning trace panels.

---

\subsection{3. Panel Specif\-ications}

\subsubsection{Panel 1: AutoApprover Audit Log}
\begin{lstlisting}
html = studio.get_autoapprover_panel(n=50)
\end{lstlisting}
Scrollable table: finding ID, entity pair, decision, confidence, tier reached (hard gate / SGD / LLM), timestamp. Color-coded by decision (green=APPROVE, red=REJECT, yellow=REVIEW).

\subsubsection{Panel 2: Contra\-dictionResolver Revision Queue}
\begin{lstlisting}
html = studio.get_revision_queue_panel()
\end{lstlisting}
Table of pending revision candidates: entity pair, net\_evidence\_score, resolution class, time in queue, nomination count.

\subsubsection{Panel 3: DiscoveryCalibrator Community Heatmap}
\begin{lstlisting}
fig = studio.get_calibrator_heatmap()
\end{lstlisting}
Horizontal bar chart: one bar per community, length = inverse-rate weight. Communities with weight \textgreater  2.0 highlighted in yellow (under\-explored). Communities with weight \textless  0.5 highlighted in grey (saturated). Cold-start communities (never scanned) shown in blue.

\subsubsection{Panel 4: ChemicalModulator Blood Panel}
\begin{lstlisting}
fig = studio.get_modulator_panel()
\end{lstlisting}
Five-scalar visua\-lization comparing current levels to homeostatic baselines:
\begin{itemize}
\item Horizontal bars: current value vs. baseline
\item Overdriven (\textgreater  1.5$\times$ baseline) $\to$ orange
\item Depleted (\textless  0.5$\times$ baseline) $\to$ blue
\item Normal $\to$ green

\end{itemize}
\subsubsection{Panel 5: Autonomous Loop Cycle History}
\begin{lstlisting}
html = studio.get_loop_panel()
\end{lstlisting}
Table of \texttt{CycleRecord} objects: cycle \#, timestamp, scan count, mater\-ializations, approvals, rejections, reviews, circuit breaker status, effective\_cap (Phase 82), edges\_rolled\_back (Phase 79). Status badge: RUNNING / STOPPED / TRIPPED.

\subsubsection{Panel 6: Provenance Panel (Phase 78)}
\begin{lstlisting}
stats_html, batch_fig, timeline_fig = studio.get_provenance_panel(n=20)
\end{lstlisting}

Three components:
\begin{enumerate}
\item \textbf{Summary row} (4 HTML cards): total batches, total edges mater\-ialized, total edges rolled back, active batch count.
\item \textbf{Batch bar chart}: horizontal bars per batch sorted newest-first. Green = active; red = rolled back. Truncated to \texttt{n} most recent.
\item \textbf{Cycle timeline}: dual-series chart. Per-cycle mater\-ialization count (bars, left axis) + cumulative edges mater\-ialized (dashed line, right axis). Rolled-back batches reduce the cumulative line.

\end{enumerate}
---

\subsection{4. Graceful Degradation}

Every panel method returns a safe placeholder when the required engine is not attached:

\begin{lstlisting}
# Panel 3 without calibrator attached:
fig = studio.get_calibrator_heatmap()
# → empty Figure with annotation "DiscoveryCalibrator not configured"
\end{lstlisting}

This is enforced by the \texttt{@requires\_attachment} decorator pattern, ensuring no \texttt{AttributeError} propagates to the Gradio UI even in minimal deployments.

---
\textbf{Copyright © 2026 Bryan Alexander Buchorn. All Rights Reserved.}

---
\subsection{Acknow\-ledgments}

The author gratefully ackno\-wledges the use of Claude (Anthropic) as a research assistant throughout this work. Claude assisted with mathe\-matical forma\-lization, code generation, manuscript preparation, and technical writing. All conceptual contr\-ibutions, archi\-tectural decisions, exper\-imental design, and intel\-lectual claims are solely the author's.

\textbf{Reviewed on}: May 2, 2026 for version v2.52.0
